\documentclass[11pt,a4paper]{article}
\newcommand{\intestazione}{Dimero di spin --- 11. Readout asimmetrico}
% =====================================================================
%  Preambolo comune ai documenti sul dimero (serie "che cosa abbiamo fatto")
%  Incluso con % =====================================================================
%  Preambolo comune ai documenti sul dimero (serie "che cosa abbiamo fatto")
%  Incluso con % =====================================================================
%  Preambolo comune ai documenti sul dimero (serie "che cosa abbiamo fatto")
%  Incluso con \input{preambolo_dimero} subito dopo \documentclass.
% =====================================================================
\usepackage[T1]{fontenc}
\usepackage[utf8]{inputenc}
\usepackage[main=italian,provide=*]{babel}
\usepackage{amsmath,amssymb,mathtools}
\usepackage{graphicx}
\usepackage{booktabs}
\usepackage{colortbl}
\usepackage{xcolor}
\usepackage{geometry}
\usepackage{placeins}
\usepackage{caption}
\usepackage{enumitem}
\usepackage{fancyhdr}
\usepackage{needspace}
\usepackage{tikz}
\usetikzlibrary{arrows.meta,positioning,shapes.geometric,calc,fit,backgrounds}
\usepackage{hyperref}
\hypersetup{colorlinks=true,linkcolor=blu,citecolor=blu,urlcolor=blu}

\geometry{margin=2.5cm}

% --- colori coerenti con quelli delle figure (palette Okabe-Ito)
\definecolor{blu}{HTML}{0072B2}
\definecolor{arancio}{HTML}{E69F00}
\definecolor{verdes}{HTML}{009E73}
\definecolor{rossos}{HTML}{D55E00}
\definecolor{violas}{HTML}{CC79A7}
\definecolor{grigetto}{HTML}{E8E8E8}

\captionsetup{font=small,labelfont=bf,width=0.92\textwidth}

\pagestyle{fancy}
\fancyhf{}
\fancyhead[L]{\small\itshape\intestazione}
\fancyhead[R]{\small\thepage}
\renewcommand{\headrulewidth}{0.3pt}

% --- notazione
\newcommand{\ket}[1]{\lvert #1\rangle}
\newcommand{\bra}[1]{\langle #1 \rvert}
\newcommand{\media}[1]{\langle #1 \rangle}
\newcommand{\vs}[1]{\mathbf{s}_{#1}}

% --- riquadro "in una riga": il messaggio da portare a casa
\newcommand{\messaggio}[1]{%
  \begin{center}
  \begin{minipage}{0.93\textwidth}
  \setlength{\fboxsep}{9pt}%
  \colorbox{grigetto}{\begin{minipage}{\dimexpr\textwidth-2\fboxsep}
  \small\textbf{In una riga.}\enspace #1
  \end{minipage}}
  \end{minipage}
  \end{center}
}

% --- riquadro "come lo sappiamo": la verifica numerica dietro un'affermazione
\newcommand{\verifica}[1]{%
  \begin{center}
  \begin{minipage}{0.93\textwidth}
  \small\textcolor{verdes}{\rule{2.2em}{1.1pt}}\enspace
  \textcolor{verdes}{\textbf{Come lo sappiamo.}}\enspace #1
  \end{minipage}
  \end{center}
  \medskip
}
 subito dopo \documentclass.
% =====================================================================
\usepackage[T1]{fontenc}
\usepackage[utf8]{inputenc}
\usepackage[main=italian,provide=*]{babel}
\usepackage{amsmath,amssymb,mathtools}
\usepackage{graphicx}
\usepackage{booktabs}
\usepackage{colortbl}
\usepackage{xcolor}
\usepackage{geometry}
\usepackage{placeins}
\usepackage{caption}
\usepackage{enumitem}
\usepackage{fancyhdr}
\usepackage{needspace}
\usepackage{tikz}
\usetikzlibrary{arrows.meta,positioning,shapes.geometric,calc,fit,backgrounds}
\usepackage{hyperref}
\hypersetup{colorlinks=true,linkcolor=blu,citecolor=blu,urlcolor=blu}

\geometry{margin=2.5cm}

% --- colori coerenti con quelli delle figure (palette Okabe-Ito)
\definecolor{blu}{HTML}{0072B2}
\definecolor{arancio}{HTML}{E69F00}
\definecolor{verdes}{HTML}{009E73}
\definecolor{rossos}{HTML}{D55E00}
\definecolor{violas}{HTML}{CC79A7}
\definecolor{grigetto}{HTML}{E8E8E8}

\captionsetup{font=small,labelfont=bf,width=0.92\textwidth}

\pagestyle{fancy}
\fancyhf{}
\fancyhead[L]{\small\itshape\intestazione}
\fancyhead[R]{\small\thepage}
\renewcommand{\headrulewidth}{0.3pt}

% --- notazione
\newcommand{\ket}[1]{\lvert #1\rangle}
\newcommand{\bra}[1]{\langle #1 \rvert}
\newcommand{\media}[1]{\langle #1 \rangle}
\newcommand{\vs}[1]{\mathbf{s}_{#1}}

% --- riquadro "in una riga": il messaggio da portare a casa
\newcommand{\messaggio}[1]{%
  \begin{center}
  \begin{minipage}{0.93\textwidth}
  \setlength{\fboxsep}{9pt}%
  \colorbox{grigetto}{\begin{minipage}{\dimexpr\textwidth-2\fboxsep}
  \small\textbf{In una riga.}\enspace #1
  \end{minipage}}
  \end{minipage}
  \end{center}
}

% --- riquadro "come lo sappiamo": la verifica numerica dietro un'affermazione
\newcommand{\verifica}[1]{%
  \begin{center}
  \begin{minipage}{0.93\textwidth}
  \small\textcolor{verdes}{\rule{2.2em}{1.1pt}}\enspace
  \textcolor{verdes}{\textbf{Come lo sappiamo.}}\enspace #1
  \end{minipage}
  \end{center}
  \medskip
}
 subito dopo \documentclass.
% =====================================================================
\usepackage[T1]{fontenc}
\usepackage[utf8]{inputenc}
\usepackage[main=italian,provide=*]{babel}
\usepackage{amsmath,amssymb,mathtools}
\usepackage{graphicx}
\usepackage{booktabs}
\usepackage{colortbl}
\usepackage{xcolor}
\usepackage{geometry}
\usepackage{placeins}
\usepackage{caption}
\usepackage{enumitem}
\usepackage{fancyhdr}
\usepackage{needspace}
\usepackage{tikz}
\usetikzlibrary{arrows.meta,positioning,shapes.geometric,calc,fit,backgrounds}
\usepackage{hyperref}
\hypersetup{colorlinks=true,linkcolor=blu,citecolor=blu,urlcolor=blu}

\geometry{margin=2.5cm}

% --- colori coerenti con quelli delle figure (palette Okabe-Ito)
\definecolor{blu}{HTML}{0072B2}
\definecolor{arancio}{HTML}{E69F00}
\definecolor{verdes}{HTML}{009E73}
\definecolor{rossos}{HTML}{D55E00}
\definecolor{violas}{HTML}{CC79A7}
\definecolor{grigetto}{HTML}{E8E8E8}

\captionsetup{font=small,labelfont=bf,width=0.92\textwidth}

\pagestyle{fancy}
\fancyhf{}
\fancyhead[L]{\small\itshape\intestazione}
\fancyhead[R]{\small\thepage}
\renewcommand{\headrulewidth}{0.3pt}

% --- notazione
\newcommand{\ket}[1]{\lvert #1\rangle}
\newcommand{\bra}[1]{\langle #1 \rvert}
\newcommand{\media}[1]{\langle #1 \rangle}
\newcommand{\vs}[1]{\mathbf{s}_{#1}}

% --- riquadro "in una riga": il messaggio da portare a casa
\newcommand{\messaggio}[1]{%
  \begin{center}
  \begin{minipage}{0.93\textwidth}
  \setlength{\fboxsep}{9pt}%
  \colorbox{grigetto}{\begin{minipage}{\dimexpr\textwidth-2\fboxsep}
  \small\textbf{In una riga.}\enspace #1
  \end{minipage}}
  \end{minipage}
  \end{center}
}

% --- riquadro "come lo sappiamo": la verifica numerica dietro un'affermazione
\newcommand{\verifica}[1]{%
  \begin{center}
  \begin{minipage}{0.93\textwidth}
  \small\textcolor{verdes}{\rule{2.2em}{1.1pt}}\enspace
  \textcolor{verdes}{\textbf{Come lo sappiamo.}}\enspace #1
  \end{minipage}
  \end{center}
  \medskip
}


\title{\bfseries Quando la scala moltiplicativa non basta più\\[2pt]
\large Documento 11 --- il test del relatore sul readout asimmetrico}
\author{Samuele Schiavi \\ \small Tesi triennale in Fisica, Università di Parma
--- Relatore: Prof.\ Alessandro Chiesa}
\date{}

\begin{document}
\maketitle
\thispagestyle{fancy}

\section{Una domanda lasciata esplicitamente aperta}

Il relatore, discutendo il modello di rumore, ha chiesto di partire da un
readout simmetrico e poi --- tempo permettendo, senza investirci troppo ---
provare anche un test con readout asimmetrico. Il readout simmetrico è già
stato usato ovunque nella pipeline (Documenti 5--10), con il risultato,
dimostrato analiticamente nel Documento 9, che $N^*=2$ del correlatore
dinamico è indipendente da $p_\text{readout}$ per \emph{qualunque} valore
simmetrico. Questo documento affronta la parte lasciata aperta: cosa succede
quando il readout smette di essere simmetrico.

\section{Perché non è la ripetizione di un conto già chiuso}

Nel caso simmetrico, un solo parametro $p$ descrive il readout: la
probabilità di leggere il risultato sbagliato è la stessa in entrambe le
direzioni, $0\to1$ e $1\to0$, e la correzione $\langle
Z\rangle_\text{readout}=(1-2p)\langle Z\rangle_\text{ideale}$ è un fattore
puramente moltiplicativo --- da cui l'invarianza di $N^*$ già verificata nel
Documento 9.

Un readout asimmetrico ha invece due parametri distinti,
$p_{01}=P(\text{misuro }1\mid\text{stato }\ket0)$ e
$p_{10}=P(\text{misuro }0\mid\text{stato }\ket1)$, in generale diversi.
Ripetendo la stessa derivazione --- probabilità lette come combinazione
lineare pesata dalla matrice di confusione, poi la differenza
$\langle Z\rangle=P(0)-P(1)$, poi la sostituzione $P_0=(1+Z)/2$,
$P_1=(1-Z)/2$ --- il risultato non è più una semplice riscalatura:
\begin{equation}
\langle Z\rangle_\text{readout} = \alpha\,\langle Z\rangle_\text{ideale} + \beta,
\qquad \alpha \equiv 1-p_{01}-p_{10}, \qquad \beta \equiv p_{10}-p_{01}.
\label{eq:formula-generale}
\end{equation}
Una trasformazione \textbf{affine}, non più puramente moltiplicativa: al
fattore di scala $\alpha$ (analogo diretto di $1-2p$) si aggiunge un
\textbf{offset costante} $\beta$, indipendente dallo stato misurato. Ponendo
$p_{01}=p_{10}=p$ si ottiene $\alpha=1-2p$, $\beta=0$: il caso noto è il
sottocaso $\beta=0$ della formula generale, non una coincidenza ma la verifica
che la derivazione più ampia contiene quella già fatta.

\messaggio{Il caso simmetrico non si estende automaticamente. Per il
correlatore, l'osservabile rilevante è un \textbf{modulo} nel piano complesso,
e sommare una costante prima di prendere il modulo non commuta con
l'operazione di modulo: in generale $|\alpha C(N)+\beta(1+i)|\neq
\alpha|C(N)|+\text{cost.}$ per nessuna scelta di costante. Geometricamente, il
readout simmetrico ($\beta=0$) riscala il luogo dei punti $C(N)$ nel piano
complesso rispetto all'origine --- il punto più lontano resta il più lontano.
Il readout asimmetrico ($\beta\neq0$) trasla l'intera curva di un vettore
costante: la distanza dall'origine di ciascun punto cambia in modo diverso a
seconda della sua posizione originaria, e il massimo dopo la traslazione può
non essere più quello di prima. L'invarianza di $N^*$ è garantita solo se
$\beta=0$; per $\beta\neq0$ non c'è nessuna garanzia analoga, va verificato
numericamente --- è esattamente questo che rende la domanda del relatore
genuinamente aperta, non un conto già deciso in anticipo.}

\section{Da dove viene l'asimmetria, fisicamente}

Il meccanismo dominante dell'errore di lettura nei qubit superconduttori è il
rilassamento ($T_1$) durante la finestra di misura: un qubit preparato in
$\ket1$ ha una probabilità non nulla di decadere a $\ket0$ prima che la misura
sia completata, mentre non esiste un meccanismo altrettanto diretto per cui un
qubit in $\ket0$ decada spontaneamente a $\ket1$ --- violerebbe la
conservazione dell'energia in assenza di eccitazione esterna. Questo produce
un'asimmetria strutturale $p_{10}\gtrsim p_{01}$, osservata sistematicamente
nelle proprietà di backend pubblicate da IBM, tipicamente di un fattore
$2$--$4\times$ su molti dispositivi Heron e precedenti.

\section{Una ricerca che non ha trovato tutto quello che cercava}

La fonte già citata in tesi per la calibrazione di \texttt{ibm\_torino}
(Stenger \emph{et al.}, arXiv:2504.15187) riporta solo la mediana simmetrica
$p_\text{readout}=2.3\times10^{-2}$, non la scomposizione in $p_{01}$ e
$p_{10}$ separati. È stata condotta una ricerca esplicita di uno split
realmente registrato per \texttt{ibm\_torino}: un solo paper (AbuGhanem,
arXiv:2410.00916) presenta una tabella dedicata con i valori per-qubit, ma
pubblicata come immagine non estraibile in forma testuale citabile con i
numeri esatti. Nessun'altra fonte citabile riporta lo split per questo chip
specifico.

La ricerca ha però trovato split reali, citabili, su altri chip IBM
comparabili per tecnologia e ordine di grandezza dell'errore:

\begin{center}
\renewcommand{\arraystretch}{1.3}
\small
\begin{tabular}{@{}llccc@{}}
\toprule
Fonte & Chip & $p_{01}$ & $p_{10}$ & Rapporto $p_{10}/p_{01}$\\
\midrule
arXiv:2203.14756 & IBM-Quito, Q0 & $0.0132$ & $0.0468$ & $3.5\times$\\
Documentazione IBM & qubit d'esempio & $0.0295$ & $0.0908$ & $3.1\times$\\
arXiv:2504.12575 & ibmq\_montreal, Q0 & $0.008$ & $0.016$ & $2.0\times$\\
\bottomrule
\end{tabular}
\end{center}

A questi si aggiunge una conferma qualitativa dalla documentazione IBM stessa:
l'errore \texttt{prob\_meas0\_prep1} è di norma il fattore dominante
nell'errore di lettura, proprio a causa del rilassamento durante la misura ---
$p_{10}>p_{01}$ è la norma su hardware superconduttore reale, non
un'eccezione.

\messaggio{Onestà dichiarata, non aggirata: non esiste, con le fonti già
citate in tesi, uno split asimmetrico realmente registrato per
\texttt{ibm\_torino}. I valori usati nel test che segue sono
\textbf{illustrativi} --- un rapporto rappresentativo tratto da calibrazioni
IBM reali su chip comparabili, stesso trattamento già riservato a
$\varepsilon_{2q}$ nel Documento 5.}

\section{I valori scelti, e perché}

Rapporto $p_{10}:p_{01}=3$ --- nel mezzo dell'intervallo $2$--$3.5\times$
trovato in letteratura --- alla \textbf{stessa media} già in uso,
$p_\text{readout}=2.3\times10^{-2}$: $p_{01}=0.0115$, $p_{10}=0.0345$. Tenere
fissa la media isola l'effetto \emph{dell'asimmetria in sé}, senza mescolarlo
con un cambiamento della scala complessiva del rumore di lettura, già
scansionato nel caso simmetrico al Documento 9.

\section{Verifica prima di fidarsi del risultato}

Tre controlli, nell'ordine in cui è naturale farli --- tutti su dati
genuinamente campionati a $2\times10^5$ shot (Documento 8), non su una
formula:

\begin{itemize}[leftmargin=1.4em,itemsep=0.2em]
\item \textbf{Limite di rumore nullo} ($p_{01}=p_{10}=0$): coincide col ramo
originale (readout spento) entro il rumore statistico atteso, su tutti gli
$N$ testati.
\item \textbf{Limite simmetrico} ($p_{01}=p_{10}=p$): riproduce $N^*=2$ e
valori $C(N)$ compatibili, entro il rumore statistico, con quelli già
registrati nel Documento 8 --- il caso noto contenuto nella formula
generale, come atteso dall'Eq.~\eqref{eq:formula-generale}, questa volta
verificato su un campionamento indipendente (seme del simulatore diverso),
non ricalcolato con la stessa formula.
\item \textbf{Risultato centrale}: $N^*=2$ non si sposta con lo split
illustrativo, né con rapporti fino a $50\times$ (Sez.~\ref{sec:stress}).
\end{itemize}

\section{Risultati: il correlatore, simmetrico contro asimmetrico}

\begin{figure}[htbp]
\centering
\includegraphics[width=0.72\textwidth]{figure/fig_correlatore_N_asimmetrico}
\caption{Modulo del correlatore rumoroso $|C_{11}^{yz}(t{=}2,N)|$ in funzione
di $N$, readout genuinamente campionato: simmetrico ($p_{01}=p_{10}=0.023$)
contro asimmetrico ($p_{01}=0.0115$, $p_{10}=0.0345$, stessa media). Il
massimo resta a $N^*=2$ in entrambi i casi. Lo split asimmetrico alza
$|C(N)|$ su \textbf{tutti} gli $N$ testati, non solo a $N$ grande: l'offset
$\beta(1+i)$ si somma qui in modo costruttivo con $C_\text{ideale}(N)$ su
tutto l'intervallo --- non un effetto casuale, ma dipendente dalla fase
specifica di $C_\text{ideale}(N)$ per questo problema.}
\end{figure}

\FloatBarrier

\begin{center}
\renewcommand{\arraystretch}{1.2}
\begin{tabular}{@{}lcc@{}}
\toprule
Configurazione di readout & $N^*$ & $|C(N^*)|$\\
\midrule
Simmetrico ($p_{01}=p_{10}=0.023$) & $2$ & $0.5869$\\
Asimmetrico illustrativo ($p_{01}=0.0115$, $p_{10}=0.0345$) & $2$ & $0.6195$\\
\bottomrule
\end{tabular}
\end{center}

\section{Uno stress test oltre il range realistico}
\label{sec:stress}

Per capire quanto il risultato $N^*=2$ sia vicino a un limite o robusto, il
test è stato ripetuto a media fissa con rapporti ben oltre l'intervallo
realistico trovato in letteratura ($2$--$3.5\times$):

\begin{center}
\renewcommand{\arraystretch}{1.15}
\small
\begin{tabular}{@{}cccc@{}}
\toprule
Rapporto $p_{10}{:}p_{01}$ & $p_{01}$ & $p_{10}$ & $N^*$\\
\midrule
$1$ (simmetrico) & $0.0230$ & $0.0230$ & $2$\\
$3$ (illustrativo) & $0.0115$ & $0.0345$ & $2$\\
$5$ & $0.0077$ & $0.0383$ & $2$\\
$10$ & $0.0042$ & $0.0418$ & $2$\\
$20$ & $0.0022$ & $0.0438$ & $2$\\
$50$ & $0.0009$ & $0.0451$ & $2$\\
\bottomrule
\end{tabular}
\end{center}

\begin{figure}[htbp]
\centering
\includegraphics[width=0.62\textwidth]{figure/fig_Nstar_vs_rapporto}
\caption{$N^*$ in funzione del rapporto $p_{10}{:}p_{01}$ (scala log), a
media fissa $2.3\times10^{-2}$. La linea resta piatta a $N^*=2$ per quattro
ordini di grandezza di rapporto, ben oltre il range realistico
($2$--$3.5\times$).}
\end{figure}

\FloatBarrier

\messaggio{$N^*$ risulta empiricamente molto più robusto di quanto la teoria
garantisse. La teoria non prometteva l'invarianza per $\beta\neq0$ --- il
conto era genuinamente aperto. Numericamente, per questo correlatore del
dimero a questo punto di lavoro, $N^*=2$ è invariante non solo nel range
realistico ma fino a un rapporto $50\times$, ben oltre qualunque asimmetria
osservabile su hardware reale.}

\section{Perché non ha ceduto nemmeno a rapporti estremi}

Guardando la figura del correlatore, l'offset $\beta(1+i)$ sposta la curva
$C_\text{ideale}(N)$ in una direzione che, per questo correlatore, si
allontana dall'origine più agli $N$ già vicini al picco che agli $N$ lontani
da esso: l'effetto netto è accentuare il picco esistente piuttosto che
spostarlo. Questo non era garantito dalla teoria --- che ammette anche il caso
opposto, un offset che appiattisce o sposta il massimo --- ma è coerente con
essa: entrambe le possibilità erano geometricamente plausibili prima di
sostituire numeri.

\section{Considerazioni}

\begin{itemize}[leftmargin=1.5em,itemsep=3pt]
\item \textbf{La domanda era genuinamente aperta, la risposta è empirica.} A
differenza del caso simmetrico, dove l'invarianza di $N^*$ è un teorema, qui
la teoria lascia il risultato indeterminato: l'esito --- $N^*$ invariante, con
ampio margine --- è un fatto numerico specifico di questo correlatore e di
questo punto di lavoro, non una conseguenza necessaria della struttura del
problema. Non è detto che valga per un correlatore diverso, un $t$ diverso, o
un altro punto di lavoro.
\item \textbf{Nessun valore reale registrato per \texttt{ibm\_torino} esiste
in forma citabile}, nonostante la ricerca esplicita. Dichiarato con la stessa
onestà con cui si dichiarano le altre scelte di modellazione della tesi: il
risultato numerico qui presentato usa un rapporto rappresentativo, non lo
split vero del chip citato altrove per gli altri parametri.
\item \textbf{Criterio di proporzionalità rispettato.} Il relatore ha chiesto
di non investire tempo eccessivo su questo punto opzionale: la teoria più
questo test numerico chiudono la domanda con lo stesso livello di rigore degli
altri stadi della pipeline, senza aprire un'indagine sproporzionata rispetto
alla natura facoltativa della richiesta.
\end{itemize}

\messaggio{Il test opzionale del relatore è chiuso: l'asimmetria del readout
non sposta $N^*=2$ (correlatori dinamici), né al punto illustrativo (rapporto
$3\times$, stessa media) né a rapporti molto più estremi ($50\times$) del
range osservato su hardware IBM reale. Questo risultato non era garantito
dalla teoria --- a differenza del caso simmetrico, dove l'invarianza è
dimostrata analiticamente --- ed è quindi un fatto verificato, non assunto.}

\section{Dove chiude la Parte 2 sul dimero}

Con questo documento si chiude, sul dimero, l'intera pipeline di rumore:
cinque stadi (Documenti 5--9) più due estensioni facoltative (Documenti 10 e
11), tutte verificate numericamente, non solo argomentate. Il prossimo passo
del progetto è replicare la stessa pipeline sul sistema a tre spin, con la
topologia ad anello --- una scelta discussa e confermata con il relatore dopo
aver confrontato anello e catena aperta sul caso a sistema chiuso.

\end{document}
